\chapter{Linear Algebra Review --- Solutions}

\noindent\textit{Source: Justin Wyss-Gallifent, \textit{Math 401/431: Linear Algebra Review}.
These exercises cover the prerequisite material pointed to in SP1.}

% -------------------------------------------------------------------------
\section*{Section 7: Eigenvalues and Eigenvectors}

\begin{exercise}{7.1}
Find the characteristic polynomials and eigenvalues of:\\
(a) $A=\begin{pmatrix}4&-5\\1&-2\end{pmatrix}$\quad
(b) $A=\begin{pmatrix}2&-1\\1&4\end{pmatrix}$\quad
(c) $A=\begin{pmatrix}-1&-2\\2&-1\end{pmatrix}$
\end{exercise}
\begin{solution}
\textbf{(a)} $\det(\lambda I-A)=(\lambda-4)(\lambda+2)+5=\lambda^2-2\lambda-3=(\lambda-3)(\lambda+1)$.
\[
\boxed{\lambda_1=3,\quad\lambda_2=-1.}
\]
Eigenvector for $\lambda_1=3$: $(3I-A)=\begin{pmatrix}-1&5\\-1&5\end{pmatrix}\to v_1=\begin{pmatrix}5\\1\end{pmatrix}$.

Eigenvector for $\lambda_2=-1$: $(-I-A)=\begin{pmatrix}-5&5\\-1&1\end{pmatrix}\to v_2=\begin{pmatrix}1\\1\end{pmatrix}$.

\textbf{(b)} $\det(\lambda I-A)=(\lambda-2)(\lambda-4)+1=\lambda^2-6\lambda+9=(\lambda-3)^2$.
\[
\boxed{\lambda=3\ (\text{algebraic multiplicity }2).}
\]
$(3I-A)=\begin{pmatrix}1&1\\-1&-1\end{pmatrix}\to$ eigenspace spanned by $\begin{pmatrix}1\\-1\end{pmatrix}$ (geometric multiplicity 1; not diagonalizable over $\mathbb{R}$).

\textbf{(c)} $\det(\lambda I-A)=(\lambda+1)^2+4=\lambda^2+2\lambda+5$.
\[
\lambda=\frac{-2\pm\sqrt{4-20}}{2}=-1\pm 2i.
\]
\[
\boxed{\lambda_{1,2}=-1\pm 2i.}
\]
(No real eigenvalues; the matrix represents a rotation-plus-scaling in $\mathbb{R}^2$.)
\end{solution}

\clearpage
% -------------------------------------------------------------------------
\section*{Section 9: Linear Transformations}

\begin{exercise}{9.1}
Determine whether each transformation is linear. If not, give a counterexample.\\
(a) $T(x_1,x_2)=(x_1+2x_2,\,x_1,\,-3x_2)$\quad
(b) $T(x_1,x_2,x_3)=(-x_1,\,x_1+x_2)$\quad
(c) $T(x_1,x_2)=(x_1+2,\,3x_2)$
\end{exercise}
\begin{solution}
\textbf{(a) Linear.}
\[
T(\alpha u+\beta v)
=\bigl(\alpha u_1+\beta v_1+2(\alpha u_2+\beta v_2),\;\alpha u_1+\beta v_1,\;-3(\alpha u_2+\beta v_2)\bigr)
=\alpha T(u)+\beta T(v).
\]
Representing matrix: $A=\begin{pmatrix}1&2\\1&0\\0&-3\end{pmatrix}$.

\textbf{(b) Linear.}
\[
T(\alpha u+\beta v)=\bigl(-(\alpha u_1+\beta v_1),\;(\alpha u_1+\beta v_1)+(\alpha u_2+\beta v_2)\bigr)
=\alpha T(u)+\beta T(v).
\]
Representing matrix: $A=\begin{pmatrix}-1&0&0\\1&1&0\end{pmatrix}$.

\textbf{(c) Not linear.}

By Theorem 9.0.1, any linear transformation must satisfy $T(\mathbf{0})=\mathbf{0}$.
\[
T(0,0)=(0+2,\;3\cdot 0)=(2,0)\neq(0,0).
\]
So $T$ is not linear. Explicit violation: $T(2\cdot(1,0))=T(2,0)=(4,0)$ but $2T(1,0)=2(3,0)=(6,0)\neq(4,0)$. $\square$
\end{solution}

\begin{exercise}{9.3}
Use the matrix method (Theorem 9.0.2) to check if $T(x_1,x_2)=(5x_1+7x_2,\;5x_1-x_2)$ is linear.
\end{exercise}
\begin{solution}
\textbf{Step 1: Assume linearity and find candidate matrix.}
\[
T(e_1)=T(1,0)=(5,5),\quad T(e_2)=T(0,1)=(7,-1).
\]
Candidate matrix: $M=\begin{pmatrix}5&7\\5&-1\end{pmatrix}$.

\textbf{Step 2: Verify $T(x_1,x_2)=Mx$.}
\[
M\begin{pmatrix}x_1\\x_2\end{pmatrix}=\begin{pmatrix}5x_1+7x_2\\5x_1-x_2\end{pmatrix}=T(x_1,x_2).\quad\checkmark
\]
\textbf{Conclusion: $T$ is linear.} The representing matrix is $M=\begin{pmatrix}5&7\\5&-1\end{pmatrix}$.
\end{solution}

\begin{exercise}{9.4}
Use the matrix method to check if $T(x_1,x_2)=(x_1^2\,x_2,\;x_1-x_2)$ is linear.
\end{exercise}
\begin{solution}
\textbf{Step 1: Assume linearity and find candidate matrix.}
\[
T(e_1)=T(1,0)=(1^2\cdot0,\;1-0)=(0,1),\quad T(e_2)=T(0,1)=(0^2\cdot1,\;0-1)=(0,-1).
\]
Candidate matrix: $M=\begin{pmatrix}0&0\\1&-1\end{pmatrix}$.

\textbf{Step 2: Verify $T(x_1,x_2)=Mx$.}
\[
M\begin{pmatrix}x_1\\x_2\end{pmatrix}=\begin{pmatrix}0\\x_1-x_2\end{pmatrix}.
\]
But $T(x_1,x_2)=(x_1^2 x_2,\,x_1-x_2)$. The first component is $x_1^2 x_2$, not $0$.

\textbf{Contradiction.} We have a specific counterexample: let $x_1=2,x_2=1$:
\[
T(2,1)=(4\cdot1,\;2-1)=(4,1),\quad M\begin{pmatrix}2\\1\end{pmatrix}=(0,1)\neq(4,1).
\]
\textbf{Conclusion: $T$ is not linear.} $\square$

(Alternative direct proof: $T(2e_1)=T(2,0)=(0,2)$ but $2T(e_1)=2(0,1)=(0,2)$ --- this particular check passes. But $T(e_1+e_2)=T(1,1)=(1,0)$ while $T(e_1)+T(e_2)=(0,1)+(0,-1)=(0,0)\neq(1,0)$.)
\end{solution}
